% FORMLÆRE: MATHEMATICA — Ei formell grunnlegging
% Iver Raknes Finne, AHO 2026
% Søstertekst til FORMLÆRE (traktatform i prosa)
% Spraak: Nynorsk

\documentclass[11pt, a4paper]{article}

\usepackage[utf8]{inputenc}
\usepackage[T1]{fontenc}
\usepackage[nynorsk]{babel}
\usepackage{mathpazo}
\usepackage{amsmath, amssymb, amsthm}
\usepackage{hyperref}
\usepackage{setspace}
\usepackage{geometry}
\usepackage[round]{natbib}
\usepackage{csquotes}
\usepackage{enumitem}
\usepackage{booktabs}
\usepackage{array}
\usepackage{longtable}

\geometry{margin=2.5cm, top=2.5cm, bottom=2.5cm}
\setstretch{1.15}

\hypersetup{
  colorlinks=true,
  linkcolor=blue!60!black,
  citecolor=green!50!black,
  urlcolor=blue!70!black
}

% --- Theorem environments ---
\theoremstyle{definition}
\newtheorem{defn}{Definisjon}[section]
\newtheorem{postulat}[defn]{Postulat}

\theoremstyle{plain}
\newtheorem{teorem}[defn]{Teorem}
\newtheorem{prop}[defn]{Proposisjon}

\theoremstyle{remark}
\newtheorem*{merknad}{Merknad}
\newtheorem*{doeme}{D\o me}
\newtheorem*{falsifisering}{Falsifiseringsvilk\aa r}
\newtheorem*{prediksjon}{Prediksjon}

% --- Custom commands ---
\newcommand{\MC}{\mathbf{M}_C}
\newcommand{\x}{\mathbf{x}}
\newcommand{\y}{\mathbf{y}}
\newcommand{\svec}{\mathbf{s}}
\newcommand{\Reals}{\mathbb{R}}
\newcommand{\formlaere}{\textsc{Forml\ae re}}

\title{%
  \vspace{-1cm}
  {\huge\textbf{FORML\AE RE: MATHEMATICA}}\\[0.5em]
  {\large\textit{Ei formell grunnlegging}}
}
\author{
  Iver Raknes Finne\\
  \texttt{iver.raknes.finne@stud.aho.no}\\
  Arkitektur- og designh\o gskolen i Oslo\\
  2026
}
\date{}

\begin{document}
\maketitle
\thispagestyle{empty}

\begin{quote}
\itshape
Denne teksten er s\o sterteksten til \formlaere. Der \formlaere{} legg ut rammeverket i prosa, byggjer \textsc{Mathematica} opp det same rammeverket i formell notasjon, fr\aa{} punkt til vektor til funksjonale. Kvar definisjon korresponderer med ein proposisjon i hovudteksten; krysstilvisingane er eksplisitte.
\end{quote}

\newpage
\tableofcontents
\newpage

% ============================================================
% I  ROM
% ============================================================
\section{Rom}
\label{sec:rom}

\begin{quote}
\itshape Korresponderer med \formlaere{} 1.1, 1.2, 1.21--1.23
\end{quote}

\begin{defn}[Attributtvektor]
\label{def:attributtvektor}
Lat ein gjenstand $x$ vere skildra av $n$ m\aa lbare fysiske eigenskapar.
Attributtvektoren til $x$ er
\[
  \x = (x_1, x_2, \ldots, x_n) \in \Reals^n
\]
der kvar komponent $x_i$ svarar til ein eigenskap: h\o gd, breidd, djupn, masse, hardheit, overflateruheit, fargeb\o lgjelengd, osb. Mengda av eigenskapar er ikkje gjeven a priori; ho vert bestemt av analyseform\aa let og kan utvidast.
\end{defn}

\begin{merknad}
Val av eigenskapar definerer ein \emph{representasjon}. Ulike representasjonar gjev ulike rom. Modellen krev berre at representasjonen er konsistent innanfor ein analyse, ikkje at det finst \'ein kanonisk representasjon.
\end{merknad}

\begin{defn}[Formrom]
\label{def:formrom}
Lat $C$ vere ein funksjonell klasse (t.d.\ ``stolar,'' ``kniver,'' ``bustadhus'').
Formrommet til $C$ er
\[
  \MC \subseteq \Reals^n
\]
mengda av alle logisk moglege attributtvektorar for klassen. $\MC$ er ikkje heile $\Reals^n$; han er avgrensa av fysiske lovar og funksjonskrav.
\end{defn}

\begin{defn}[Regiontypologi]
\label{def:regiontypologi}
Formrommet $\MC$ har tre typar regionar:
\begin{enumerate}[label=(\roman*)]
  \item \textbf{Busett region} $B = \{\,\x \in \MC : \x \text{ er realisert i minst \'ein faktisk gjenstand}\,\}$
  \item \textbf{Open region} $O = \{\,\x \in \MC : \x \text{ er fysisk mogleg men urealisert}\,\}$
  \item \textbf{Forboden region} $F = \Reals^n \setminus \MC$ \quad (fysisk umoglege former)
\end{enumerate}
slik at $\MC = B \cup O$ og $B \cap O = \varnothing$.
\end{defn}

\begin{merknad}
Grensa $\partial\MC$ mellom $O$ og $F$ er ein funksjon av tilgjengeleg teknologi $\tau$.
Me skriv $\mathbf{M}_C(\tau)$ n\aa r me vil gjere teknologiavhenget eksplisitt.
Formrommet \emph{veks} med teknologien: om $\tau_1 \subset \tau_2$, s\aa{}
$\mathbf{M}_C(\tau_1) \subseteq \mathbf{M}_C(\tau_2)$.
\end{merknad}

\begin{prop}
\label{prop:forklaringskrevjande}
Sidan $|B| > 1$ for kvar ikkje-triviell klasse $C$ (\formlaere{} 1.3),
er det eit forklaringskrevjande faktum at ein gjenstand okkuperer \'ein posisjon
$\x \in B$ og ikkje ei anna $\y \in B$.
Forklaringa av dette er emnet for alt som f\o lgjer.
\end{prop}


% ============================================================
% II  KREFTER
% ============================================================
\section{Krefter}
\label{sec:krefter}

\begin{quote}
\itshape Korresponderer med \formlaere{} 2.1--2.43
\end{quote}

\begin{defn}[Seleksjonstrykk]
\label{def:seleksjonstrykk}
Eit seleksjonstrykk er ein funksjon
\[
  s_i : \MC \to \Reals
\]
som tilordnar kvar posisjon $\x$ ein reell verdi $s_i(\x)$ som uttrykkjer kor godt $\x$ tilfredsstiller det $i$-te vilk\aa ret. H\o gare verdi er betre.
\end{defn}

\begin{doeme}
For ein stol med funksjonell klasse $C$ = ``sittem\o bel'':
\begin{itemize}
  \item $s_1(\x)$ = materialtilgang: kor lett tilgjengeleg materialkombinasjonen er
  \item $s_2(\x)$ = produksjonskostnad (negativt forteikn, sidan l\aa gare er betre)
  \item $s_3(\x)$ = brukskrav: kor godt stolen st\o r kroppen
  \item $s_4(\x)$ = kulturell aksept: kor godt forma samsvarar med omgjevnadens forventningar
  \item $s_5(\x)$ = materiell affordanse: kor lett forma let seg realisere i det gjevne materialet
  \item[] \ldots{} opp til $s_k(\x)$ for $k$ samtidige seleksjonstrykk
\end{itemize}
\end{doeme}

\begin{defn}[Seleksjonstrykkvektoren]
\label{def:seleksjonsvektor}
Settet av alle samtidige seleksjonstrykk definerer ein vektorfunksjon
\[
  \svec : \MC \to \Reals^k, \qquad
  \svec(\x) = \bigl(s_1(\x),\, s_2(\x),\, \ldots,\, s_k(\x)\bigr)
\]
\end{defn}

\begin{postulat}[Fleire uavhengige trykk]
\label{post:fleiretrykk}
For kvar funksjonell klasse $C$ gjeld $k \ge 2$, og det finst minst to trykk
$s_i, s_j$ som er statistisk uavhengige over $\MC$:
\[
  \operatorname{Cov}(s_i, s_j) \neq 0
  \quad\text{i minst \'ein region, men } s_i \text{ er ikkje ein funksjon av } s_j.
\]
\end{postulat}

\begin{falsifisering}
Postulatet fell om det finst ein klasse $C$ der \'eitt einaste $s_i$ er tilstrekkeleg til \aa{} forklare all observert variasjon i $B$.
\end{falsifisering}

\begin{teorem}[Variasjon under konstant funksjon]
\label{teo:variasjon}
Lat $s_1$ representere funksjonskravet. Innanfor ein funksjonell klasse $C$ er $s_1$
per definisjon konstant: $s_1(\x) = s_1(\y)$ for alle $\x, \y \in \MC$.
Sidan $k \ge 2$ (postulat~\ref{post:fleiretrykk}), finst det minst eitt anna trykk
$s_2$ som varierer over $\MC$. Difor bestemmer ikkje $s_1$ posisjonen aaleine,
og formvariasjon under konstant funksjon er det forventa resultatet.
\end{teorem}

\begin{proof}
Lat $\x, \y \in B$ med $s_1(\x) = s_1(\y)$. Sidan $s_2$ er uavhengig av $s_1$,
kan $s_2(\x) \neq s_2(\y)$. Difor $\x \neq \y$ sj\o lv om funksjonen er den same.
\end{proof}

\begin{defn}[Funksjonell avgrensing]
\label{def:funksjonell-avgrensing}
Funksjonen definerer klassen og dermed formrommet. Lat $F : \Reals^n \to \{0, 1\}$
vere indikatorfunksjonen for funksjonskravet. D\aa{} er
\[
  \MC = \bigl\{\,\x \in \Reals^n : F(\x) = 1\,\bigr\}
\]
Funksjonen avgrensar $\MC$. Innanfor $\MC$ er ho taus.
\end{defn}

\begin{postulat}[Motstridande retningar]
\label{post:motstridande}
For minst eitt par $(s_i, s_j)$ finst det ein region $R \subseteq \MC$ der
\[
  \nabla s_i(\x) \cdot \nabla s_j(\x) < 0
  \qquad \text{for } \x \in R
\]
det vil seie: gradientane peikar i motstridande retningar.
\AA{} forbetre \'ein eigenskap forverrar den andre.
\end{postulat}

\begin{falsifisering}
Postulatet fell om alle gradientpar samsvarar overalt i $\MC$.
\end{falsifisering}

\begin{teorem}[Kompromiss]
\label{teo:kompromiss}
Av postulata~\ref{post:fleiretrykk} og~\ref{post:motstridande} f\o lgjer det at
kvar realisert form $\x \in B$ er eit kompromiss: ein posisjon der dei motstridande
trykka er balanserte, ikkje eliminerte.
\end{teorem}

\begin{proof}
Sidan $\nabla s_i \cdot \nabla s_j < 0$ i minst \'ein region, finst det ingen posisjon
$\x \in \MC$ der alle trykk samstundes er maksimerte.
Kvar realisert $\x \in B$ representerer difor ein avveging.
\end{proof}

\begin{teorem}[Fleire gyldige kompromiss]
\label{teo:fleire-kompromiss}
Sidan $|B| > 1$ (proposisjon~\ref{prop:forklaringskrevjande}) og kvar $\x \in B$ er
eit kompromiss (teorem~\ref{teo:kompromiss}), finst det fleire gyldige kompromiss.
Skilnaden mellom to former under same funksjon er ikkje ein feil, men eit uttrykk for
at kreftene kan balanserast p\aa{} meir enn \'ein m\aa te.
\end{teorem}


% ============================================================
% III  LANDSKAP
% ============================================================
\section{Landskap}
\label{sec:landskap}

\begin{quote}
\itshape Korresponderer med \formlaere{} 3.1--3.42
\end{quote}

\begin{defn}[Tilpassingsfunksjon]
\label{def:tilpassingsfunksjon}
Tilpassingsfunksjonen er ein aggregering av seleksjonstrykka:
\[
  f : \MC \to \Reals, \qquad
  f(\x) = \Phi\bigl(s_1(\x),\, s_2(\x),\, \ldots,\, s_k(\x)\bigr)
\]
der $\Phi : \Reals^k \to \Reals$ er ein aggregeringsfunksjon som vektar dei ulike trykka.
Valet av $\Phi$ er ikkje unikt; ulike vektingar gjev ulike landskap.
\end{defn}

\begin{merknad}
I den enklaste modellen er $\Phi$ line\ae r: $f(\x) = \sum_i w_i\, s_i(\x)$ der
$w_i \ge 0$ er vekter. Men $\Phi$ kan vere ikkje-line\ae r, t.d.\ om to trykk interagerer.
\end{merknad}

\begin{defn}[Tilpassingslandskap]
\label{def:landskap}
Tilpassingslandskapet er grafen til $f$ over $\MC$:
\[
  L = \bigl\{\,(\x,\, f(\x)) : \x \in \MC\,\bigr\} \subset \Reals^{n+1}
\]
Lokale maksimum av $f$ er \textbf{haugar}. Lokale minimum er \textbf{dalar}.
Salar mellom haugar er \textbf{ryggar}.
\end{defn}

\begin{teorem}[Fleire haugar]
\label{teo:fleirehaugar}
Sidan $k \ge 2$ uavhengige trykk verkar (postulat~\ref{post:fleiretrykk}), og sidan
minst eitt par har motstridande gradientar (postulat~\ref{post:motstridande}),
har tilpassingsfunksjonen $f$ generisk fleire lokale maksimum.
\end{teorem}

\begin{proof}
Lat $s_i$ ha eit lokalt maksimum ved $\mathbf{a}$ og $s_j$ ha eit lokalt maksimum
ved $\mathbf{b} \neq \mathbf{a}$, og lat $\nabla s_i \cdot \nabla s_j < 0$ i regionen
mellom dei. For vektinga $w$ der $w_i \gg w_j$ er det eit lokalt maksimum n\ae r
$\mathbf{a}$; for vektinga $w'$ der $w_j \gg w_i$ er det eit lokalt maksimum n\ae r
$\mathbf{b}$. Sidan vektingane varierer kontinuerleg, har landskapet minst to haugar.
\end{proof}

\begin{merknad}
I NK-landskapsteori \citep{kauffman1993} aukar talet lokale maksimum med epistasen $K$.
I formrommet svarar $K$ til graden av interaksjon mellom seleksjonstrykka.
\end{merknad}

\begin{defn}[Stil]
\label{def:stil}
Ein stil $S$ er ei samanhengande delmengd av $B$ rundt eit lokalt maksimum $\x^*$ av $f$:
\[
  S = \bigl\{\,\x \in B : \|\x - \x^*\| < \varepsilon
  \;\text{og}\; f(\x) > f(\x^*) - \delta\,\bigr\}
\]
for eigna $\varepsilon, \delta > 0$.
Ein stil er ein attraksjonsbasseng i tilpassingslandskapet.
\end{defn}

\begin{teorem}
\label{teo:stilar}
Av teorem~\ref{teo:fleirehaugar} f\o lgjer det at det finst fleire stilar.
Av definisjon~\ref{def:tilpassingsfunksjon} og~\ref{def:stil} f\o lgjer det at kvar stil
er fullt ut bestemt av seleksjonstrykka og vektinga; stilen har ingen eigne friheitsgrader.
\end{teorem}

\begin{defn}[Vandring]
\label{def:vandring}
Ei formgjevingshistorie er ein stig i formrommet:
\[
  \gamma : [0, T] \to \MC, \qquad
  \gamma(0) = \x_0, \qquad
  \gamma(T) = \x_T
\]
der kvart steg er lokalt: $\|\gamma(t+\mathrm{d}t) - \gamma(t)\| < \varepsilon$ for alle $t$.
Stigen er stiavhengig: den tilgjengelege regionen ved tid $t$ avheng av $\gamma([0, t])$.
\end{defn}


% ============================================================
% IV  DYNAMIKK
% ============================================================
\section{Dynamikk}
\label{sec:dynamikk}

\begin{quote}
\itshape Korresponderer med \formlaere{} 4.1--4.6
\end{quote}

\begin{postulat}[Dynamisk landskap]
\label{post:dynamisk}
Seleksjonstrykka er tidavhengige:
\[
  s_i(\x, t) : \MC \times \Reals \to \Reals
\]
og difor er tilpassingslandskapet eit tidavhengig funksjonale:
\[
  f(\x, t) = \Phi\bigl(s_1(\x, t),\, \ldots,\, s_k(\x, t)\bigr)
\]
\end{postulat}

\begin{falsifisering}
Postulatet fell om tilpassingslandskapet for ein klasse kan visast \aa{} vere topologisk
uendra over ein periode der den observerte fordelinga av former endrar seg p\aa{} ein m\aa te
som ikkje kan forklarast av stokastisk drift.
\end{falsifisering}

\begin{teorem}[Punktert likevekt]
\label{teo:punktert}
Av teorem~\ref{teo:fleirehaugar} og postulat~\ref{post:dynamisk} f\o lgjer det at
formhistoria viser punktert likevekt.
\end{teorem}

\begin{proof}
Former samlar seg rundt haugar (teorem~\ref{teo:fleirehaugar}).
Haugane flyttar seg med landskapet (postulat~\ref{post:dynamisk}).
Lat $\x^*(t)$ vere posisjonen til eit lokalt maksimum ved tid $t$.
S\aa{} lenge $\partial\x^*/\partial t$ er liten, f\o lgjer formene gradvis med (stabilitet).
N\aa r $\partial f/\partial t$ er stort nok til at haugen forsvinn
(d.v.s.\ $\partial^2 f/\partial\x^2$ skiftar forteikn), kollapsar klyngja,
og formene m\aa{} finne nye haugar (omvelting).
\end{proof}

\begin{prediksjon}
Shannon-entropien $H'(t) = -\sum p_i \ln p_i$ over materialkategoriar i eit datasett
b\o r vise plat\aa{} avbrotne av br\aa{} stigingar.
STOLAR-datasettet ($n \approx 2\,300$, 1280--2024) gjev:
$H' = 2{,}67$ p\aa{} 1600-talet; $H' = 3{,}51$ p\aa{} 1900-talet.
Diskontinuitetane svarar til teknologiske omveltingar
(kolonial import ca.\ 1600, industrialisering ca.\ 1875).
\end{prediksjon}

\begin{defn}[Landskapsminne]
\label{def:landskapsminne}
Kvar realisert form $\x \in B$ etterlet seg eit spor som modifiserer
seleksjonstrykka for neste form:
\[
  s_i(\y, t+\mathrm{d}t) = s_i(\y, t) + \eta_i(\y, \x)
\]
der $\eta_i$ er ein ``impaktfunksjon'' som uttrykkjer korleis den realiserte forma $\x$
endrar det $i$-te trykket for framtidige former $\y$. Landskapet har minne.
\end{defn}


% ============================================================
% V  MATERIALE
% ============================================================
\section{Materiale}
\label{sec:materiale}

\begin{quote}
\itshape Korresponderer med \formlaere{} 5.1--5.52
\end{quote}

\begin{postulat}[Geometrisk signatur]
\label{post:signatur}
Kvart materiale $m$ induserer ei ikkje-uniform sannsynlegheitsfordeling over
underrommet av geometriar $G \subset \Reals^n$:
\[
  P(\x \mid m) \neq \mathrm{Uniform}(G)
  \qquad \text{for alle } m
\]
Materialet trekkjer forma mot visse konfigurasjonar og motset seg andre.
\end{postulat}

\begin{falsifisering}
Postulatet fell om same geometriske distribusjon kan observerast uavhengig av materiale,
innanfor same funksjonelle klasse og same sett av \o vrige seleksjonstrykk.
\end{falsifisering}

\noindent\textbf{Formelt:}
Lat $m_1, m_2$ vere to ulike materialar.
Postulat~\ref{post:signatur} hevdar at
$P(\x \mid m_1, C, \svec_{-m}) \neq P(\x \mid m_2, C, \svec_{-m})$
der $\svec_{-m}$ er alle seleksjonstrykk utanom materialet.

\begin{teorem}[Informasjonsteoremet]
\label{teo:informasjon}
Innanfor ein funksjonell klasse $C$ ber materialet meir gjensidig informasjon med
realisert geometri enn funksjonen gjer:
\[
  I(G;\, M) > I(G;\, \mathrm{Func}) = 0
\]
\end{teorem}

\begin{proof}
Funksjonen er per definisjon konstant innanfor $C$.
Ein konstant variabel har null entropi: $H(\mathrm{Func} \mid C) = 0$.
Difor $I(G;\, \mathrm{Func}) = H(G) - H(G \mid \mathrm{Func}) = 0$.
Etter postulat~\ref{post:signatur} er $P(G \mid M)$ ikkje-uniform,
s\aa{} $H(G \mid M) < H(G)$, og difor
$I(G;\, M) = H(G) - H(G \mid M) > 0 > I(G;\, \mathrm{Func})$.
\end{proof}

\begin{prediksjon}
I STOLAR-datasettet skal $I(\text{geometri};\, \text{materiale}) > 0$ innanfor
stolklassen. Dette kan estimerast med $k$-n\ae raste-nabo-estimatorar for gjensidig
informasjon \citep{kraskov2004}.
\end{prediksjon}


% ============================================================
% VI  NAVIGASJON
% ============================================================
\section{Navigasjon}
\label{sec:navigasjon}

\begin{quote}
\itshape Korresponderer med \formlaere{} 6.1--6.6
\end{quote}

\begin{defn}[Navigator]
\label{def:navigator}
Ein navigator er eit trippel $N = (G, \mu, \alpha)$ der:
\begin{enumerate}[label=(\alph*)]
  \item $G \subseteq \MC$ er ei mengd av m\aa ltilstandar (ein region navigatoren styrer mot)
  \item $\mu : \MC \to \Reals_{\ge 0}$ er ein avstandsfunksjon:
        $\mu(\x) = d(\x, G)$ for ein eigna metrikk $d$
  \item $\alpha : \MC \to T\MC$ er eit justeringsfelt: ein vektorfunksjon som dreg
        systemet mot $G$, slik at
        \[
          \langle \alpha(\x),\, -\nabla\mu(\x) \rangle > 0
          \qquad \text{for } \mu(\x) > 0
        \]
\end{enumerate}
Det vil seie: $\alpha$ peikar i retning av avtakande avstand til $G$.
Navigatoren justerer seg mot m\aa let.
\end{defn}

\begin{merknad}
Definisjonen krev korkje medvit, intensjon eller nervesystem.
Ho krev berre at tre vilk\aa r er oppfylte.
Ein termostat er ein navigator.
Ein gradient-descent-algoritme er ein navigator.
Ein kvar kjemisk reaksjon som driv mot likevekt er ein navigator.
\end{merknad}

\begin{defn}[Grenseflate]
\label{def:grenseflate}
Grenseflata til ein navigator $N$ er den maksimale regionen i rom-tid der $N$
kan ha m\aa l:
\[
  \partial N = \bigl\{\,(r, t) \in \Reals^3 \times \Reals :
  N \text{ kan representere og arbeide mot tilstandar i } (r, t)\,\bigr\}
\]
Dette er \textbf{den kognitive lyskjegla} til $N$ (jf.\ Levin, 2019).
Storleiken p\aa{} $\partial N$ avgjer kva problem $N$ kan handsame.
\end{defn}

\medskip
\begin{center}
\begin{tabular}{@{}llll@{}}
\toprule
\textbf{Grenseflate} & \textbf{Romsleg skala} & \textbf{Tidsskala} & \textbf{Eksempel} \\
\midrule
$\mu$m, ms      & Intracellul\ae r & Augeblikkeleg & Ionekanalregulering \\
mm, minutt       & Cellul\ae r      & Kort          & Kjemotaksis \\
cm, timar        & Vev              & Medium        & S\aa rtilheling \\
m, \aa r         & Organisme        & Lang          & Anatomisk homeostase \\
km, ti\aa r      & Samfunn          & Generasjonell & Handverkstradisjon \\
Kontinent, hundre\aa r & Sivilisasjon & Historisk  & Arkitektonisk stil \\
Globalt, tusen\aa r    & Biosystemet  & Evolusjon\ae r & Artsmorfologi \\
\bottomrule
\end{tabular}
\end{center}
\medskip

\begin{defn}[Intelligens]
\label{def:intelligens}
Graden av intelligens til ein navigator $N$ er proporsjonal med hans evne til
\aa{} unng\aa{} lokale optimum i tilpassingslandskapet:
\[
  I(N) \propto \max\bigl\{\,d\bigl(\gamma_N(t),\, \x_{\mathrm{lokal}}\bigr) :
  \gamma_N \text{ konvergerer mot } \x_{\mathrm{global}}\,\bigr\}
\]
der $\x_{\mathrm{lokal}}$ er eit lokalt maksimum og $\x_{\mathrm{global}}$ er eit betre
globalt alternativ. Intelligens er evna til \aa{} f\o lgje omvegar som endar betre enn
den rette linja. Formelt er det evna til \aa{} krysse energibarrierar i landskapet.
\end{defn}

\begin{merknad}
Denne definisjonen er skala-uavhengig. Ho gjeld for eit enzym som unng\aa r ein lokal
minima i konformasjonsrom, for ein handverkar som omg\ae r ein umogleg geometri via ein
alternativ konstruksjonsmetode, og for ein sivilisasjon som investerer i infrastruktur
(kortsiktig kostnad) for langsiktig vinst.
\end{merknad}


% ============================================================
% VII  SUBSTRAT-UAVHENGIGHEIT
% ============================================================
\section{Substrat-uavhengigheit}
\label{sec:substrat}

\begin{quote}
\itshape Korresponderer med \formlaere{} 7.1--7.32
\end{quote}

\begin{postulat}[Substrat-uavhengigheit]
\label{post:substrat}
Definisjonen av navigator (definisjon~\ref{def:navigator}) krev ingen bestemt fysisk
substrat. Vilk\aa ra (a), (b), (c) kan oppfyllast i eit vilk\aa rleg fysisk system som har
tilstandar, kan registrere avstand fr\aa{} ein m\aa ltilstand, og kan justere seg.
\end{postulat}

\begin{falsifisering}
Postulatet fell om det kan visast at minst eitt av vilk\aa ra (a), (b), (c) krev ein
eigenskap som berre organisk materie har.
\end{falsifisering}

\begin{teorem}
\label{teo:substrat}
Av definisjon~\ref{def:navigator} og postulat~\ref{post:substrat} f\o lgjer det at
navigasjon er ein eigenskap ved systemets funksjonelle organisering, ikkje ved materialet
det er laga av. Kvar dei tre vilk\aa ra er oppfylte, finst navigasjon.
\end{teorem}

\noindent
Me viser no korleis navigator-trippelet $(G, \mu, \alpha)$ vert realisert i fem
ulike substrat:

\bigskip
\noindent\textbf{VII.3.1\quad Diffusjonsmodell.}

\medskip
\begin{center}
\begin{tabular}{@{}l >{\raggedright\arraybackslash}p{8cm}@{}}
\toprule
\textbf{Komponent} & \textbf{Realisering} \\
\midrule
Formrom $\mathbf{M}$         & Latentrom $\Reals^d$ \\
M\aa ltilstand $G$            & Kondisjonert distribusjon $p(\x_0 \mid c)$ \\
Avstandsfunksjon $\mu$        & Score $\nabla_{\x} \log p(\x_t)$ \\
Justeringsfelt $\alpha$       & Stokastisk differensiallikning (reverse SDE) \\
Seleksjonstrykk $\svec$       & Treningsdata \\
Grenseflate $\partial N$      & Parameterrom og kondisjoneringsrom \\
\bottomrule
\end{tabular}
\end{center}

\begin{merknad}
Moduskollaps under sterk kondisjonering er same fenomen som konvergens mot ein haug
(\formlaere{} 3.41): eitt seleksjonstrykk dominerer, og variasjonen snevrar inn.
\end{merknad}

\bigskip
\noindent\textbf{VII.3.2\quad Embryo / Morfogenese.}

\medskip
\begin{center}
\begin{tabular}{@{}l >{\raggedright\arraybackslash}p{8cm}@{}}
\toprule
\textbf{Komponent} & \textbf{Realisering} \\
\midrule
Formrom $\mathbf{M}$         & Morforom: rom av moglege anatomiar \\
M\aa ltilstand $G$            & Anatomisk settform (kodifisert i bioelektrisk m\o nster) \\
Avstandsfunksjon $\mu$        & Avvik mellom noverande og m\aa l-bioelektrisk m\o nster \\
Justeringsfelt $\alpha$       & Cellevandring, differensiering, apoptose, proliferasjon \\
Seleksjonstrykk $\svec$       & Genetiske, bioelektriske, mekaniske, kjemiske signal \\
Grenseflate $\partial N$      & Heil organisme (i regenererande artar); organprimordium (i ikkje-regenererande) \\
\bottomrule
\end{tabular}
\end{center}

\begin{merknad}
Planaria som regenererer hovud etter amputasjon navigerer tilbake til same posisjon i
morforom uavhengig av kor dei vert kutta.
Det bioelektriske m\o nsteret kan overskrivast utan \aa{} endre genomet (Levin, 2022):
m\aa ltilstanden $G$ er redigerbar, ikkje hardkoda.
\end{merknad}

\bigskip
\noindent\textbf{VII.3.3\quad Handverk (M\o belsnekkar).}

\medskip
\begin{center}
\begin{tabular}{@{}l >{\raggedright\arraybackslash}p{8cm}@{}}
\toprule
\textbf{Komponent} & \textbf{Realisering} \\
\midrule
Formrom $\mathbf{M}$         & Formrom for stolklassen $\MC$ \\
M\aa ltilstand $G$            & Mental representasjon + teikningar + tradisjon \\
Avstandsfunksjon $\mu$        & Taktil og visuell tilbakemelding \\
Justeringsfelt $\alpha$       & Kvart kutt, kvar b\o ying, kvar limsaum \\
Seleksjonstrykk $\svec$       & Materiale, \o konomi, tradisjon, kroppens bruk, estetikk \\
Grenseflate $\partial N$      & Verkstaden og levetida til handverkaren \\
\bottomrule
\end{tabular}
\end{center}

\begin{merknad}
Materialet navigerer samstundes mot sin eigen geometriske signatur
(postulat~\ref{post:signatur}). Den ferdige forma oppst\aa r \emph{mellom} handverkaren
og materialet. Ingen av dei kontrollerer utfallet aaleine.
\end{merknad}

\bigskip
\noindent\textbf{VII.3.4\quad Evolusjon\ae r prosess.}

\medskip
\begin{center}
\begin{tabular}{@{}l >{\raggedright\arraybackslash}p{8cm}@{}}
\toprule
\textbf{Komponent} & \textbf{Realisering} \\
\midrule
Formrom $\mathbf{M}$         & Genotyp-fenotyp-rom \\
M\aa ltilstand $G$            & H\o gare fitness (implisitt, ikkje representert) \\
Avstandsfunksjon $\mu$        & Differensiell reproduksjon \\
Justeringsfelt $\alpha$       & Seleksjon, mutasjon, drift, rekombinasjon \\
Seleksjonstrykk $\svec$       & Milj\o, konkurranse, seksuell seleksjon \\
Grenseflate $\partial N$      & Populasjonen og dens generasjonstid \\
\bottomrule
\end{tabular}
\end{center}

\begin{merknad}
Ingen einskild organisme har oversikt. Populasjonen vandrar i landskapet utan ein
sentralisert plan. Det er distribuert navigasjon (\formlaere{} 8) i sin reinaste form.
\end{merknad}

\bigskip
\noindent\textbf{VII.3.5\quad Spr\aa kmodell (LLM).}

\medskip
\begin{center}
\begin{tabular}{@{}l >{\raggedright\arraybackslash}p{8cm}@{}}
\toprule
\textbf{Komponent} & \textbf{Realisering} \\
\midrule
Formrom $\mathbf{M}$         & Tokenrom $V^T$ for sekvensar av lengd $T$ \\
M\aa ltilstand $G$            & Kondisjonert fordeling $P(x_{t+1} \mid x_{1:t}, \text{prompt})$ \\
Avstandsfunksjon $\mu$        & Tap-funksjonen (kryssentropi) \\
Justeringsfelt $\alpha$       & Autoregresjon: kvart generert token er eit steg i vandringa \\
Seleksjonstrykk $\svec$       & Treningskorpus, RLHF, promptkontekst \\
Grenseflate $\partial N$      & Kontekstvinduet og parameterrommet \\
\bottomrule
\end{tabular}
\end{center}

\begin{merknad}
Kvar generert sekvens er ein \emph{veg} gjennom tokenrommet, ikkje ein posisjon.
Sekvensen som vert realisert er den som gjer verknaden stasjon\ae r (\formlaere{} 9):
ho minimerer kryssentropien langs vegen, ikkje berre ved endepunktet.
Temperaturparameteren kontrollerer spreiinga (\formlaere{} 9.41).
\end{merknad}

\bigskip

\begin{defn}[Overtalbarheitskontinuumet]
\label{def:overtalbarheit}
Lat $P(N)$ vere den optimale styringsstrategien for navigator $N$.
Navigatorar ligg p\aa{} eit kontinuum:
\begin{description}[font=\normalfont\itshape, leftmargin=2.5cm, labelwidth=1.8cm]
  \item[Klasse A:] $P(N)$ = maskinvare-modifikasjon
    (ingen settpunkt; alt er hardkoda)
  \item[Klasse B:] $P(N)$ = settpunkt-omskriving
    (m\aa ltilstand er redigerbar, men ikkje l\ae rbar)
  \item[Klasse C:] $P(N)$ = trening med bel\o ning/straff
    (systemet l\ae rer av erfaring)
  \item[Klasse D:] $P(N)$ = kommunikasjon av grunnar
    (systemet responderer p\aa{} argument)
\end{description}
For kvart steg opp treng ein \emph{mindre kunnskap om systemets indre} og \emph{meir
kommunikasjon med systemet som heilskap}. Plasseringa av ein navigator p\aa{} dette
kontinuumet avgjer kva slags interaksjon som er optimal.
\end{defn}

\begin{merknad}
Kontinuumet er sj\o lv ein dimensjon i eit \emph{metaformrom}: rommet av moglege
navigatorar. \formlaere{} er ein modell som opererer i dette metarommet.
\end{merknad}

\end{document}
